%!TEX root = TQFTnotes.tex
%%%%%%%%%%%%%%%%%%%%%%%%%%%%%%
\chapter{3-category of tensor categories}\label{c:3cat-tens}
%References: \cite{dsps-balanced}, \cite{dsps-dualizable},
% \cite{etingof-fusion}.

In the previous chapters, we have given two constructions of an extended
2-dimensional TQFT with values in the 2-category $\Alg$ of algebras and
bimodules: one construction was based on the work of Schommer-Pries, describing
the 2-category $\Bordor_2$ by generators and relations, and the other was based
on triangulations (or more general PL cell decompositions) of a 2-manifold. In
both cases, such a theory is determined by $A=Z(\bullet)$ which must be a
separable symmetric Frobenius algebra. This is a special case of the general
cobordism hypothesis, according to which extended framed $n$-dimensional TQFTs
with values in $\calC$ are described by fully dualizable objects in $\calC$, and
oriented extended TQFTs are described by $\SO(n)$ homotopy fixed points in
$\calC^{\fd}$. In particular, for $n=2$, $\calC=\Alg$, fully dualizable objects
are separable algebras, and $\SO(2)$ homotopy fixed points are
separable symmetric Frobenius algebras.

In this chapter, we start exploring the 3d analog of these constructions. As we
will show, in this case one possible replacement for the 2-category $\Alg$ is
the 3-category $\Tens$ of finite tensor categories, and the fully dualizable
objects in it are exactly semisimple tensor categories, i.e., fusion categories;
$\SO(3)$ homotopy fixed points are (conjecturally) spherical fusion categories.
The corresponding 3-dimensional oriented extended TQFT is known as Turaev--Viro
theory.

In this chapter, we will define and study the 3-category $\Tens$, following
the works \cite{dsps-balanced}, \cite{dsps-dualizable}; the construction of
Turaev--Viro theory will be given in a later chapter.

To simplify the exposition, throughout this chapter we assume that $\kk$ is an
algebraically closed field of characteristic zero. As usual, we will skip most
proofs, referring the reader to original papers.

\section{Review of 2-category $\Alg$}\label{s:alg-review}
For comparison, we recall the main steps in the construction of the 2-category $\Alg$:
\begin{itemize}
  \item We start with the category $\vect$ of $\kk$-vector spaces and define
    the tensor product functor $\otimes\colon\vect\times\vect \to \vect$
  \item Next, we define an algebra $A$ as an object $A\in\vect$ together with
    a morphism $m\colon A\otimes A\to A$ and an element $1\in A$ satisfying
    associativity and unit axioms. We also define the notion of morphisms between
    algebras.
  \item We define the notion of left and right $A$-module and $A$-$B$-bimodules,
  together with module homomorphisms.
  \item Finally, we define the notion of tensor product over an algebra: given right
  $A$-module $M$ and left $A$-module $N$, we define $M\otimes_A N$
\end{itemize}

Let us now proceed to defining categorical analogs of all of the above.

%%%%%%%%%%%%%%%%%%%%%%%%%%%%%%
\section{Finite $\kk$-linear categories}\label{s:finite-categories}
Our first step is defining a 2-category which will replace the category $\vect$.
For that, recall that we have defined the notion of a locally finite
$\kk$-linear abelian category as a $\kk$-linear abelian category where all
objects have finite length and all $\Hom$ spaces are finite-dimensional, see
\seref{s:deligne-product}. Such categories form a 2-category $\Linlf$; this
2-category should be considered a natural categorical analog of the category of
vector spaces.

Let us now consider a more restrictive set of requirements.

\begin{definition}\label{d:finite-linear}
  A $\kk$-linear category $\calA$ is called \emph{finite}
  \index{category!finite} if it satisfies the
  following properties:
  \begin{enumerate}
    \item For any $A,B\in \calA$, the vector space $\Hom_\calA(A,B)$ is
      finite-dimensional.

    \item Any object $A\in \calA$ has finite length.
    \item $\calA$ has enough projectives: any $A\in \calA$ is a quotient of
     a projective object $P\in \calA$.
    \item There are only finitely many isomorphism classes of simple objects.
  \end{enumerate}
\end{definition}
Note that we do not assume that $\calA$ is semisimple.

\begin{example}\label{x:finite-cat-alg}
  Let $A$ be a finite-dimensional algebra over $\kk$. Then the category $A\mod$
  of finite-dimensional $A$-modules is a finite $\kk$-linear category. Moreover,
  it is easy to show that conversely, every finite $\kk$-linear category is
  equivalent to the category $A\mod$ for some finite-dimensional algebra $A$
  (see \cite[Section~1.8]{etingof-fusion}).
\end{example}


Unless stated otherwise, all functors between $\kk$-linear categories will
be assumed to be additive, $\kk$-linear, and right exact.


\begin{definition}\label{d:linf}
  We denote by $\Linf$ the 2-category of finite $\kk$-linear abelian
  categories together with additive and $\kk$-linear right exact
  functors between them and functorial morphisms.
\end{definition}
The so-defined 2-category $\Linf$ is the categorical analog of the category of
finite-dimensional vector spaces.



\begin{remark}
  One might notice that we are not quite consistent: in the $n=2$ case, we started
  with category of all vector spaces, not necessarily finite-dimensional ones,
  so our 2-category $\Alg$ also included infinite-dimensional algebras ---
  which, however, were irrelevant for constructing TQFTs: any fully dualizable
  (i.e., separable) algebra is automatically finite-dimensional. In the $n=3$ case,
  we start with finite linear categories from the beginning.

  The reason is mostly technical; many of the constructions below work for any
  locally finite category (i.e., a category only satisfying the first 2
  conditions of \deref{d:finite-linear}). I expect that one can in fact define a
  version of the (not yet constructed) 3-category $\Tens$ starting with
  locally finite linear categories rather than finite; however, it does create
  some technical difficulties, and since we expect that the only categories we
  will need for extended 3d TQFTs will be finite anyway, we can save some
  trouble and only work with finite linear categories from the beginning.
  \end{remark}

%%%%%%%%%%%%%%%%%%%%%%%%%%%%
\section{Finite tensor categories}\label{s:tensor-categories}
Our next step is defining a categorical analog of the notion of associative
algebra. The obvious candidate is the notion of a monoidal category as
discussed in \chref{c:categories}. However, in order to make some future
constructions possible, we will add some additional conditions.

\begin{definition}\label{d:finite-tensor}
  A \emph{finite tensor category}\index{finite tensor category} is a finite $\kk$-linear category $\calC$
  together with a structure of a monoidal category such that
  \begin{enumerate}
    \item The tensor product functor $\otimes\colon
     \calC\times \calC\to \calC$ is $\kk$-bilinear.
    \item $\Hom(\one, \one)=\kk$.
    \item $\calC$ is rigid: every $X\in \calC$ has left and right dual (see
      \deref{d:rigid}).
  \end{enumerate}
\end{definition}
(Note that we do not give a definition of (not necessarily finite) tensor
category since the only tensor categories we will need are the finite ones.)

Condition (1) is a natural compatibility between the monoidal structure and
structure of $\kk$-linear category; condition (2) is a harmless condition which
is there just for convenience --- almost everything below can be repeated,
possibly with minor changes, without that assumption (in fact,
\cite{dsps-balanced} do not include this condition). But the rigidity condition
(3) is absolutely crucial: without it, most of the theory below would not work.


\begin{lemma}\label{l:tensor-biexact}
  Let $\calC$ be a finite tensor category. Then
  \begin{enumerate}
    \item Functor $\otimes \colon \calC\times \calC\to \calC$ is bi-exact.
    \item Object $\one\in \calC$ is simple.
  \end{enumerate}
\end{lemma}
A proof of this lemma can be found in \cite[Proposition
4.2.1, Theorem 4.3.8]{etingof-fusion}.

As an immediate corollary, we see that tensor product can be considered as a
functor $\calC\boxtimes \calC\to \calC$, where $\boxtimes$ is Deligne's product
of $\kk$-linear categories, see \seref{s:deligne-product} (in the same way as
for algebras, multiplication can be considered as either a bilinear map $A\times
A\to A$ or as a linear map $A\otimes A\to A$).

\begin{lemma}\label{l:tensor-boxtimes}
  Let $\calC$, $\calD$ be finite tensor categories. Then category
  $\calC\boxtimes \calD$ has a natural structure of a finite tensor category,
  with tensor product functor given by
  $$
  (c_1 \boxtimes d_1)\otimes (c_2\boxtimes d_2)
    = (c_1\otimes c_2)\boxtimes (d_1\otimes d_2).
  $$
\end{lemma}
This lemma is an analog of the statement that tensor product of algebras has a
natural structure of an algebra. We skip the proof.

Below are some examples of finite tensor categories.

\begin{example}\label{x:vecf-tensor}
  The category $\vecf$ of finite-dimensional vector spaces over $\kk$ is a
  finite tensor category. This theory is the categorical analog of the trivial
  algebra $A=\kk$.
\end{example}


\begin{example}\label{x:repg-tensor}
  Let $G$ be a finite group. The category $\Rep G$ of finite-dimensional
  representations of $G$ is a finite tensor category.
\end{example}
Note that if we replaced finite group by a compact group, the resulting category
would not be finite.

\begin{example}\label{x:hopf-tensor}
  Let $H$ be a finite-dimensional Hopf algebra. Then category of
  finite-dimensional $H$-modules is a finite tensor category.
\end{example}

\begin{example}\label{x:vectg-tensor}
  Let $G$ be a finite group. Let $\vect_G$ be the category of finite-dimensional
  $G$-graded vector spaces as defined in \exref{x:vecg}.
  Then $\vect_G$ is a finite tensor category.
\end{example}

As for algebras, for finite tensor categories we have the notion of ``opposite''
multiplication: if $\calC$ is a finite tensor category, we denote by
$\calC^\mop$ the same category but with opposite tensor product:
$A\otimes^{\mop} B=B \otimes A$. (Notation $\mop$ is an abbreviation for
``monoidally opposite''; it should not be confused with $\calC^\op$, which is
obtained from $\calC$ by reversing all arrows.)


%%%%%%%%%%%%%%%%%%%%%%%%%%%%%%
\section{Module categories}\label{s:module-categories}
Our next step is defining a categorical analog of a module over an algebra.

\begin{definition}\label{d:module-cat}
  Let $\calC$ be a finite tensor category. A finite left \emph{$\calC$-module
  category}\index{module category} is a finite $\kk$-linear category $\calM$
  together with a $\kk$-bilinear action functor
  $$
  \otimes\colon \calC \times \calM \to \calM
  $$
  which is right exact in $\calC$ argument, and with functorial isomorphisms
  \begin{align*}
  &(c_1\otimes c_2)\otimes m \simeq c_1 \otimes (c_2\otimes m)\\
  & \one \otimes m \simeq m
  \end{align*}
  which satisfy the obvious pentagon and unit axioms.
\end{definition}
Note that in many references (e.g., \cite{dsps-balanced}), condition that
$\otimes$ is right exact is not part of the definition. However, this condition
is necessary for most of what follows, so if we do not include it as part of the
definition, we would need to add it as an additional assumption in our theorems.
\begin{theorem}\label{t:module-action-exact}
  Let $\calM$ be a finite module category over a finite tensor category $\calC$.
  Then the action functor $\calC\otimes \calM\to\calM$ is exact in
  $\calC$ and also exact in $\calM$.
\end{theorem}
We refer the reader to \cite[Corollary 2.26]{dsps-balanced} for the proof.

In a similar way, one defines the notion of a \emph{right} module category and a
notion of a $\calC$-$\calD$ bimodule category; for the latter, we also need to add
an isomorphism
$$
  (c\otimes m)\otimes d \simeq c\otimes (m\otimes d)
$$
again satisfying obvious compatibility conditions.

\begin{example}\label{x:tensor-self-bimodule}
  Any finite tensor category $\calC$ is automatically a bimodule category
  over itself.
\end{example}

\begin{example}\label{x:linear-cat-module}
  Any finite linear category has a natural structure of a module category over
  $\vecf$.
\end{example}

\begin{example}\label{x:coset-module}
  Let $G$ be a finite group and let $H\subset G$ be a subgroup. Then the
  category $\Rep H$ is a module category over $\Rep G$.
\end{example}

\begin{example}\label{x:algebra-in-cat}
  Let $A$ be an algebra in $\calC$, i.e., an object $A\in \calC$ together with
  morphisms $m\colon A\otimes A\to A$, $i\colon \one\to A$ satisfying obvious
  associativity and unit axioms (see \cite[Section 7.8]{etingof-fusion}). A
  right $A$-module is an object $M\in \calC$ together with the action morphism
  $M\otimes A\to M$, again satisfying obvious axioms. We denote by
  $\Mod_\calC(A)$ the category of right modules over $A$, with the obvious
  morphisms between them. Then the category $\Mod_\calC(A)$ is a finite module
  category over $\calC$. Moreover, it can be shown (see
  \cite[Corollary 7.10.5]{etingof-fusion}) that every finite $\calC$ module
  category can be obtained in this way. For $\calC=\vecf$, this is exactly the
  statement given in \exref{x:finite-cat-alg}.

  This construction allows one to reduce the study of module categories over
  $\calC$ (which in general are difficult to describe) to the study of modules
  over an algebra, which is a much more explicit question. This is one of the
  main tools used to prove results about module categories; proofs of most
  theorems below are based on this approach. However, since we skip most
  proofs, we will not go into more details about algebras in categories.
\end{example}


\begin{definition}\label{d:module-functor}\index{module functor}
  Let $\calM$, $\calN$ be finite left module categories over a finite tensor
  category $\calC$. A module functor $F\colon \calM\to \calN$ is a $\kk$-linear
  functor together with functorial isomorphisms
  $$
    J\colon F(c\otimes m)\isoto c\otimes F(m)
  $$
  satisfying the obvious pentagon and triangle relations. Unless explicitly
  stated otherwise, we will always assume that $F$ is right exact; we will
  denote the set of all right exact module functors $\calM\to \calN$ by
  $$
    \Fun_\calC(\calM, \calN)=\text{right exact module functors }\calM\to\calN.
  $$
\end{definition}

%\begin{remark}
% One can consider weakening of this definition: instead of requiring $J$ to be
% an isomorphism, we can just require it to be functorial morphism. This gives
%a
% definition of \emph{lax} module functor. However, for module categories over
%a
% tensor category, it turns out that this actually gives the same set of
% functors: every lax functor is automatically strong functor, see
% \cite[Lemma 2.10]{dsps-balanced}.
%\end{remark}

Given two functors $F_1, F_2\in \Fun_\calC(\calM, \calN)$, we define
$\Hom(F_1, F_2)$ to be the set of functorial morphisms $F_1\to F_2$ which agree
with isomorphisms $J_1, J_2$; we leave it to the reader to write the
compatibility condition explicitly (or find it in
\cite[Section 7.2]{etingof-fusion}). This turns $\Fun_\calC(\calM, \calN)$
into a category and gives rise to the 2-category $\calC\mod$ of (left) finite
$\calC$-module categories, together with $\calC$-module functors and functorial
morphisms.

In a similar way one defines module functors for right module categories
and bimodule categories.


\begin{theorem}\label{t:module-functors}
  Let $\calM$, $\calN$ be finite module categories over a finite tensor
   category $\calC$. Then the category $\Fun_\calC(\calM, \calN)$ is
   a finite $\kk$-linear category.

   If, in addition, $\calC, \calM_1, \calM_2$ are semisimple, then
   $\Fun_\calC(\calM, \calN)$  is also semisimple.
\end{theorem}
Proof of the first part can be found in \cite[Proposition
7.11.6]{etingof-fusion}. (Note that in the book, they did not spell out exact requirements on
the module categories; however, if you read the proof, you will see that it
works for any finite module categories).

The second part is \cite[Theorem 2.16]{eno}.

\begin{example}\label{x:module-cat-dual}
  Let $\calM$ be a $\calC$-$\calD$ bimodule category. Define
  $$
  \calM^\vee = \Fun_\calC(\calM, \calC).
  $$
  This category is naturally a finite
  $\calD$-$\calC$ bimodule category (compare with \ecref{ec:dual-module-frob}).

  In a similar way, we can define category
  $$
    {}^\vee\calM= \Fun_{\calD^\mop}(\calM, \calD)
  $$
  which is again a $\calD$-$\calC$ bimodule category.
\end{example}




\begin{lemma}\label{l:bimodule-2category}
  Let $\calC$, $\calD$ be finite tensor categories. The 2-category of finite
  $\calC$-$\calD$ bimodule categories is equivalent to the 2-category of finite
  $(\calC\boxtimes \calD^{\mop})$ module categories. In particular, 2-category
  of
  left $\calC$-module categories is equivalent to the 2-category of right
  $\calC^\mop$-module categories.
\end{lemma}



%
%\subsection{Drinfled center}
%Recall that for a bimodule $M$ over an algebra $A$, we had the notion of
%invariant:
%$$
% M^A=HH^0(M)=\{m\in M\st am=ma\}=\Hom_{}(A, M)
%$$
%where $\Hom$ is in the category of bimodules. In particular, for $M=A$, this
%gives the center of $A$: $HH^0(A)=z(A)$.
%
%An analog of this construction for a module category $\calM$ over a tensor
%category $\calC$ is the category
%$$
% HH^0(\calM)=\Fun^{re}(\calC, \calM)
%$$
%where $\Fun^{re}$ is the category of right exact functors of $\calC$-bimodule
%categories.
%
%THis can be described in more explicit form. Every such functor $F$ defines an
%object $X=F(\one)$ together with functorial isomorphisms
%$$
%c\otimes X\simeq F(c)\simeq X\otimes c
%$$
%\begin{definition}
% Let $\calM$ be a finite $\calC$-bimodule category. We define its center
% to be the category whose ojects are objects $X\in \calM$ together with
% functorial isomorphisms $\ga_c\colon c\otimes X\isoto X\otimes c$, satisfying
% hexagon axiom: FIXME.
%
% Morphisms in $z(\calM)$ are defined as morphisms in $\calM$ which are
% compatible with $\ga_c$.
%\end{definition}
%\begin{lemma}
% One has an equivalence of categories $z(\calM)\simeq \Fun^{re}(\calC,
%\calM)$.
%\end{lemma}
%
%e cexample: drinfeld center of a category. Example: $\vect_G$.


%%%%%%%%%%
\section{Balanced tensor product}\label{s:balanced-product}

So far, we have defined a notion of a finite tensor category and a notion of
module category over a finite tensor category; these notions are natural
analogs of the notion of a commutative algebra and modules over algebras. Our
next goal is defining an analog of tensor product over $A$.

Recall that given an algebra $A$, a right $A$-module $M_A$, and a left
$A$-module ${}_AN$, there are several ways to define tensor product $M
\otimes_A N$. The most straightforward is defining $M\otimes_A N$ as a suitable
quotient of $M\otimes N$. However, a better approach is defining $M\otimes_A N$
via a universal property: namely, it is uniquely determined by the condition
that for every vector space $L$, we have a functorial isomorphism
$$
\Hom_\kk(M\otimes_A N, L) = \Bal_A(M \times N, L)
$$
where $\Bal_A$ is the space of $A$-balanced maps, i.e., bilinear maps $M\times N
\to L$ satisfying condition
$$
f(ma, n) = f(m, an)
$$
for any $m\in M$, $n\in N$, $a\in A$.

As usual, universality immediately implies that such a vector space, if it
exists, is unique up to a unique isomorphism. Existence is not obvious and is
proved by an explicit construction.

Let us now define a categorical analog of this.

\begin{definition}\label{d:balanced}
  Let $\calC$ be a finite tensor category. Let
  $\calM_\calC$, ${}_\calC \calN$ be right and left finite $\calC$-module
  categories, and let $\calL$ be a $\kk$-linear category.

  A $\calC$-balanced functor $F\colon \calM\times \calN\to \calL$ is a
  $\kk$-bilinear functor $F\colon \calM\times \calN\to \calL$, right exact in
  each variable, together with a functorial isomorphism
  $$
    F(m\otimes c, n)\simeq F(m, c\otimes n), \qquad c\in \calC
  $$
  satisfying usual axioms (see \cite[Definition 3.1]{dsps-balanced}).
\end{definition}

One can also define a notion of morphisms of $\calC$-balanced functors. We
denote by $\Bal_\calC(\calM\times \calN, \calL)$ the category of
$\calC$-balanced functors $\calM\times \calN\to \calL$.

\begin{definition}\label{d:deligne-balanced}
  In the assumptions of \deref{d:balanced}, a balanced tensor product is a
  $\kk$-linear category $\calM\boxtimes_\calC \calN$ together with a
  $\calC$-balanced functor $\boxtimes_\calC\colon \calM\times \calN\to \calM
  \boxtimes_\calC \calN$ such that for all $\kk$-linear categories
  $\calL$, the functor $\boxtimes_\calC$
  induces an equivalence of categories
  $$
    \Bal_\calC(\calM\times \calN, \calL)\simeq
    \Fun(\calM\boxtimes_\calC \calN, \calL)
  $$
\end{definition}
It is easy to show that if such $\calM\boxtimes_\calC \calN$ exists, it is
unique; more precisely, the balanced tensor product is unique up to equivalence,
and this equivalence is in turn unique up to unique natural isomorphism. Said
another way, the 2-category of linear categories representing the balanced
tensor product is either contractible or empty. The hard part is proving
existence of $\calM\boxtimes_\calC \calN$.

\begin{theorem}\label{t:deligne-balanced}
  Let $\calC$ be a finite tensor category over $\kk$, and let $\calM$, $\calN$
  be right and left finite $\calC$-module categories. Then the balanced tensor
  product $\calM\boxtimes_\calC\calN$ exists and is a finite $\kk$-linear
  category. Moreover, the functor $\boxtimes_\calC\colon \calM\times\calN\to
  \calM\boxtimes_\calC\calN$ is exact in each argument.
\end{theorem}
A proof of this theorem can be found in \cite[Theorem~3.3]{dsps-balanced}.

\begin{example}\label{x:vecf-modules}
  Let $\calC=\vecf$ be the category of finite-dimensional vector spaces.
  Then the balanced tensor product functor $\boxtimes_{\vecf}$ coincides with
  Deligne product defined in \seref{s:deligne-product}.
\end{example}

\begin{example}\label{x:module-dual-category}
  Let $\calM$ be a finite left $\calC$-module category. Then one has a natural
  equivalence $\calC\boxtimes_\calC \calM\simeq \calM$ (compare with the isomorphism
  $A\otimes_A M\simeq M$ for a module $M$ over an algebra $A$).
\end{example}

\begin{lemma}\label{l:bimodule-composition}
  Let $\calA$, $\calB$, $\calC$ be finite tensor categories, and let
  $\calM$ be a finite $\calA$-$\calB$ bimodule category, and $\calN$ a finite
  $\calB$-$\calC$ bimodule category. Then the balanced tensor product
  $\calM\boxtimes_{\calB} \calN$ has a natural structure of $\calA$-$\calC$
  bimodule category.
\end{lemma}
This lemma can be found in \cite[Remark 3.7]{dsps-balanced}.

\section{Category $\Tens$}\label{s:tens-category}

\begin{theorem}\label{t:tens-exists}
  There exists a $(3, 3)$ category $\Tens$ with the following objects and
  morphisms:
  \begin{itemize}
    \item Objects: $\kk$-linear finite tensor categories
    \item 1-morphisms $\calC\to\calD$: finite linear $\calD$-$\calC$ bimodule
      categories; horizontal composition of
      1-morphisms is given by balanced Deligne product.
    \item 2-morphisms: right exact bimodule functors
    \item 3-morphisms: functorial morphisms
  \end{itemize}
  Category
  $\Tens$ has a structure of a symmetric monoidal 3-category, with respect
  to Deligne product $\boxtimes$. The unit for this monoidal structure is the
  tensor category $\vecf$.
\end{theorem}
This theorem is stated in this form in
\cite[Theorem 2.2.18]{dsps-dualizable}.

\begin{remark}
  In my opinion, the current state of this and related statements is somewhat
  unsatisfactory. \cite[Theorem 2.2.18]{dsps-dualizable} does not give a
  proof; instead, they refer the reader to \cite{jfs}, where this statement
  seems to be stated in Example~8.10. That example does not state the theorem in
  any detail but briefly mentions that one can construct such a 3-category by
  thinking about tensor categories as algebras in a (suitably defined)
  2-category of $\kk$-linear categories, and provides a reference to yet another
  100-page paper where the theory of $E_n$ algebras in the setting of infinity
  categories is discussed. I have no doubt that the result is correct, and
  can be extracted from these papers, but I would much appreciate a
  more direct proof.
\end{remark}


Note that \cite{jfs} use a different model of higher categories than we do
(namely, their approach is based on complete Segal spaces); however, since, as
we discussed in \chref{c:higher-cat}, all models of $(\infty,n)$ categories are equivalent,
this should not be a problem.
